% Statement and cited-input packet. Target proof is separately embargoed.
% Frozen II source SHA256: 74497bfe327fae25c9cdcb17334b3bf42a9793a9df7c0a4404d97568f2f760f0
% Target owner: cdp:st:C4g; tag: 01FQ

\begin{proposition}[The necessary normalization defect in the three named test families]\hypertarget{tag-01FQ}{}\label{cdp:st:C4g}


For every admissible member of the original $(3,4)$ family,
every member of the two-parameter family in \ref{cdp:st:E13}, and every
member of the Liu--Rollenske family in \ref{cdp:st:E14}, there is a reduced
full hexagonal fibre $W$ with normalization
$\eta:\mathrm dP_6\to W$ such that
\[
 \eta^*L=\OO_{\mathrm dP_6}\quad(L\in\Pic(Y)),\qquad
 H^2\bigl(\mathrm dP_6,K_{\mathrm dP_6}\otimes
                    \eta^*(K_Y^{-1}\otimes L)\bigr)=\C.
 \tag{C.TEST.1}\hypertarget{tag-012O}{}\label{cdp:eq:C.TEST.1}
\]
In these three families the singular-fibre Euler contributions are
\[
\begin{array}{c|c|c}
\text{family}&\text{reduced singular supports}&
                 \text{their Euler numbers}\\ \hline
(3,4)&W&2\\
\text{two-parameter \ref{cdp:st:E13}}&W,\ W_1,\ W_2&2,0,0\\
\text{Liu--Rollenske \ref{cdp:st:E14}}&W,\ C_{\rm cusp}&2,0 .
\end{array}
\tag{C.TEST.2}\hypertarget{tag-012P}{}\label{cdp:eq:C.TEST.2}
\]
The other critical full fibres are smooth multiple fibres,
with multiplicities $(3,4)$, $4$, and $4$, respectively.
Every one of those full multiple fibres passes the generic
ambient vanishing test \ref{cdp:st:A.3.4}, including all its nilpotent
layers. Consequently positive $c_3$ in these examples does
not require a failure of that test on a nonreduced fibre.
The reduced full fibre $W$ already supplies the fixed
normalization obstruction, and $H^2(Y,T_Y\otimes L)\ne0$
for every ambient line bundle in each family.
\end{proposition}

% Cited input cdp:st:A.3.4; statement or labelled formula only.
\begin{proposition}[The smooth-multiple assertion of Remark 7.2]\hypertarget{tag-01AO}{}\label{cdp:st:A.3.4}


Under the hypotheses of CDP20 Proposition~7.3, let $mS$, with $m\ge2$, be a full fibre whose reduced support $S$ is smooth. For every real irrational transgression parameter $r$ of \ref{cdp:st:D.1.9},
\[
 H^0\bigl(mS,(\Omega_Y^1\otimes L_r)|_{mS}\bigr)=0.
 \tag{A.15a}\hypertarget{tag-0045}{}\label{cdp:eq:A.15a}
\]
The same simultaneous vanishing on all smooth multiple fibres $(m\ge2)$ holds on the complement of a proper complex analytic subset of $\Pic(Y)$. The assertion includes the entire scheme-theoretic fibre, with every nilpotent layer.
\end{proposition}

% Cited input cdp:st:A.4.1; statement or labelled formula only.
\begin{proposition}[Both fibre groups on the specified $W$]\hypertarget{tag-01AQ}{}\label{cdp:st:A.4.1}


For the reduced normal-crossings fibre $W$ with normalization $\eta:\operatorname{dP}_6\to W$ in the specified threefold $X$, and every $L\in\Pic(X)$, put $A=(K_X\otimes L^{-1})|_W$. Then
\[
H^0(W,Q\otimes A)=0,
\qquad H^0(W,\Omega_X^1|_W\otimes A)\simeq\C.
\tag{A.15}\hypertarget{tag-0044}{}\label{cdp:eq:A.15}
\]
The same two conclusions hold with $A=B|_W$ for every $B\in\Pic(X)$.
\end{proposition}

% Cited input cdp:st:C.1.3; statement or labelled formula only.
\begin{proposition}\hypertarget{tag-01BP}{}\label{cdp:st:C.1.3}


For every $Y$ and $D$ in Statement \CompendiumRef{II}{cdp:st:C.1.2} and every $L\in\Pic(Y)$,
\[
 h^2(Y,T_Y\otimes L)\geq
 \max\left\{0,1-\frac{b_3(Y)}2\right\}.
\]
In particular, $b_3(Y)=0$ gives $H^2(Y,T_Y\otimes L)\ne0$ for
every $L$. A threefold in this class possessing a twist with vanishing
second tangent cohomology has $b_3(Y)\geq2$.
\end{proposition}

% Cited input cdp:st:C5b; statement or labelled formula only.
\begin{proposition}[Integral twists and the normalization-map exception]\hypertarget{tag-01FS}{}\label{cdp:st:C5b}


For $n\in\Z$ define
\[
 r_1(n)\equiv n\pmod3,\quad 0\le r_1(n)<3,
 \qquad r_2(n)\equiv3n\pmod4,\quad0\le r_2(n)<4,
\]
\begin{gather*}
 d(n)=\frac{n-4r_1(n)-3r_2(n)}{12},\qquad
 \epsilon_n=\begin{cases}0&3\mid n,\\1&3\nmid n.\end{cases}
\\
 u(n)=\max\{d(n)+1,0\},\qquad
 v(n)=\max\{-d(n)-1,0\}.
\end{gather*}
The coarse direct images of $T_X\otimes M_n$ are
\begin{align*}
 f_*(T_X\otimes M_n)&=\OO(d(n))\oplus\OO(d(n-1)),\\
 R^1f_*(T_X\otimes M_n)&=\OO(d(n))\oplus\OO(d(n-1))\oplus\OO(d(n+2)),\\
 R^2f_*(T_X\otimes M_n)&=\OO(d(n+6))\oplus\C_{p_1}^{\epsilon_n}.\tag{C.5}\hypertarget{tag-0096}{}\label{cdp:eq:C.5}
\end{align*}
Their global cohomology dimensions are
\begin{align*}
 h^0&=u(n)+u(n-1),\\
 h^1&=v(n)+v(n-1)+u(n)+u(n-1)+u(n+2),\\
 h^2&=v(n)+v(n-1)+v(n+2)+u(n+6)+\epsilon_n,\\
 h^3&=v(n+6).                                             \tag{C.6}\hypertarget{tag-0097}{}\label{cdp:eq:C.6}
\end{align*}
The natural map $H^2(X,T_X\otimes M_n)\to H^2(\mathrm dP_6,K_{\mathrm dP_6})=\C$ has rank one exactly for
\[
 n\in\{-6,-3,-2\}\cup\Z_{\ge0};                         \tag{C.7}\hypertarget{tag-0098}{}\label{cdp:eq:C.7}
\]
it is zero for all other integers. In particular $M_{-1}$ has tangent cohomology $(0,0,1,0)$ and a zero map to the normalization target.
\end{proposition}

% Cited input cdp:st:E13; statement or labelled formula only.
\begin{proposition}[The Kuranishi germ of the two-parameter family]\hypertarget{tag-01GK}{}\label{cdp:st:E13}


For every member of the family of
\cite[Proposition labelled \texttt{family}]{Viaclovsky2026Two}, on the parameter
bidiscs with $\mu$ sufficiently close to zero and nonzero sufficiently
small linearization scale $\lambda$, one has
\[
 (h^0,h^1,h^2,h^3)(X_{\mu,\lambda},T_X)=(1,2,2,0).
\]
Its Kuranishi germ is a smooth reduced surface, and the displayed
$(\mu,\lambda)$-family is locally miniversal. In particular its
sufficiently small deformations all have algebraic dimension one.
\end{proposition}

% Cited input cdp:st:E14; statement or labelled formula only.
\begin{proposition}[The Liu--Rollenske family and its transverse deformation]\hypertarget{tag-01GL}{}\label{cdp:st:E14}


Every member $Y(u)$ of the Liu--Rollenske family, with its scale
sufficiently small and nonzero as in its construction, satisfies
\[
 h^\bullet(Y(u),T_{Y(u)})=(1,2,2,0).
\]
Its Kuranishi germ is a smooth reduced surface. One parameter is
the given scale; a second parameter replaces its cuspidal fibre by
two rank-one nodal fibres. All sufficiently small deformations
retain a torus fibration and have algebraic dimension one.
\end{proposition}

% Cited input cdp:st:D.1.9; statement or labelled formula only.
\begin{proposition}\hypertarget{tag-01G7}{}\label{cdp:st:D.1.9}


Let $Y$ be a connected compact complex manifold with
$b_1(Y)=b_2(Y)=0$, and let $f:Y\to\PP^1$ be a proper
holomorphic surjection with connected fibres. There is a class
$\beta\in H^1(Y,\OO_Y)$ such that, for every irrational real number
$r$, the line bundle
\[
L_r=\exp(2\pi i\,r\beta)
\]
has infinite order on every scheme-theoretic fibre of $f$.
Every connected component of $\Pic(Y)$ likewise contains a real
one-parameter family with this property at every irrational real
parameter.
No local-freeness assumption on $R^1f_*\OO_Y$ or constant-dimension
assumption on fibre cohomology is required.
\end{proposition}

